\section{Introdução}

A calculadora científica constitui uma ferramenta de apoio relevante
para estudantes de cursos da área de Ciências Exatas, especialmente em
disciplinas que envolvem a realização e a interpretação de cálculos
matemáticos. No estudo de Estatística, sua utilização pode auxiliar na
execução de cálculos repetitivos e na obtenção de resultados numéricos,
permitindo que o estudante concentre maior atenção na compreensão dos
procedimentos e na interpretação dos resultados. Apesar disso, muitos
estudantes ingressam no ensino superior sem familiaridade com os recursos
e as funcionalidades desse tipo de equipamento, o que pode dificultar sua
utilização adequada durante as atividades acadêmicas.

A utilização de recursos tecnológicos como instrumentos de apoio ao
ensino e à aprendizagem da Matemática tem sido discutida no campo da
Educação Matemática. Borba e Penteado \cite{borba2019} destacam que a
incorporação de recursos tecnológicos ao processo educacional está
relacionada às formas pelas quais professores e estudantes interagem com
o conhecimento. De maneira complementar, Pillon et al. \cite{pillon2020}
identificam, em revisão integrativa, diferentes recursos e abordagens
digitais empregados no ensino-aprendizagem de Matemática. A elaboração
de materiais voltados especificamente ao uso de calculadoras também
encontra precedentes: Becker et al. \cite{becker2020} apresentam um
minicurso de calculadora científica desenvolvido no contexto da formação
em Matemática no Instituto Federal do Rio Grande do Sul (IFRS), o que
reforça a pertinência de iniciativas dessa natureza.

Nesse contexto, este manual apresenta uma introdução ao uso da
calculadora científica Casio FX-82MS \cite{casio2004}, com foco em sua
aplicação no estudo de Estatística e destinado principalmente a
estudantes recém-ingressados em cursos da área de Tecnologia da
Informação. A escolha desse modelo está
relacionada à sua disponibilidade no mercado, à sua proposta voltada ao
público estudantil e à presença de recursos matemáticos e estatísticos
\cite{neubert2026}. As fórmulas, notações e organização dos exemplos
seguem a obra \textit{Estatística Fácil}, de Crespo \cite{crespo2009},
mantendo consistência com a disciplina de Probabilidade e Estatística
do curso Bacharelado em Ciência da Computação (IFPR -- Campus Pinhais) 
e estabelecendo relação direta entre a formulação matemática
dos problemas e sua aplicação prática por meio da calculadora.

O objetivo é promover a autonomia dos estudantes no uso da calculadora,
apresentando seus recursos de maneira gradual e contextualizada, com
exemplos que relacionam os procedimentos realizados na calculadora aos
conceitos matemáticos envolvidos. A calculadora é apresentada como
ferramenta de apoio à resolução dos problemas, e não como substituta da
compreensão dos conceitos estatísticos.
